\documentclass[12pt]{article}
%% BUGS:
%% - The \begin{problem} ... \end{problem} is only used in exams??

\usepackage{url, graphicx}
\PassOptionsToPackage{normalem}{ulem} % magic
\usepackage{href-ul}

% page layout
\usepackage[letterpaper, textheight=9.5in, textwidth=5.5in]{geometry}
\usepackage{setspace}
\setstretch{1.08}
\pagestyle{plain}
\sloppy\sloppypar\raggedbottom\frenchspacing\thispagestyle{empty}

% problem set header
\newcommand{\psname}{Problem Set}
\newcommand{\startps}[1]{\section*{Computational Physics --- \psname~{#1}}}

% problem formatting
\usepackage{amsthm}
\theoremstyle{definition}
\newcommand{\problemname}{Problem}
\newtheorem{problem}{\problemname}
\newcommand{\probref}[1]{\problemname~\ref{#1}}
\theoremstyle{remark}
\newcommand{\notename}{Note}
\newtheorem{note}{\notename}[problem]
\newcommand{\noteref}[1]{\textit{\notename}~\ref{#1}}
\newcounter{subpart}[problem]
\renewcommand{\thesubpart}{\alph{subpart}}
\newcommand{\subpart}{\par\addvspace{0.75ex}\refstepcounter{subpart}\textsl{(\thesubpart)}~}
\newcommand{\subpartref}[1]{\textsl{(\ref{#1})}}

% words
\newcommand{\foreign}[1]{\textsl{#1}}
\newcommand{\vs}{\foreign{vs}}

% code
\newcommand{\code}[1]{\texttt{\detokenize{#1}}}
\usepackage{xcolor}
% Define custom colors for code highlighting
\definecolor{codegreen}{rgb}{0,0.6,0}
\definecolor{codegray}{rgb}{0.5,0.5,0.5}
\definecolor{codepurple}{rgb}{0.58,0,0.82}
\definecolor{backcolour}{rgb}{0.98,0.98,0.99}
\usepackage{listings}
\lstdefinestyle{pythonstyle}{
    backgroundcolor=\color{backcolour},   
    commentstyle=\color{codegreen},
    keywordstyle=\color{magenta},
    numberstyle=\tiny\color{codegray},
    stringstyle=\color{codepurple},
    tabsize=4
}
\lstset{
    basicstyle=\small\ttfamily,
    columns=fixed,
    breaklines=false,
    keepspaces=true,
    showspaces=false,
    showstringspaces=false,
    showtabs=false,
    showlines=*,
    style=pythonstyle,
    language=Python,
    literate={~}{{\url{~}}}1
}

% math
\usepackage{amsmath}
\newcommand{\dd}{\mathrm{d}}
\newcommand{\e}{\mathrm{e}}

% primary units
\newcommand{\rad}{\mathrm{rad}}
\newcommand{\kg}{\mathrm{kg}}
\newcommand{\m}{\mathrm{m}}
\newcommand{\s}{\mathrm{s}}

% secondary units
\renewcommand{\deg}{\mathrm{deg}}
\newcommand{\km}{\mathrm{km}}
\newcommand{\cm}{\mathrm{cm}}
\newcommand{\mm}{\mathrm{mm}}
\newcommand{\ft}{\mathrm{ft}}
\newcommand{\mi}{\mathrm{mi}}
\newcommand{\AU}{\mathrm{a.u.}}
\newcommand{\au}{\mathrm{a.u.}}
\newcommand{\ns}{\mathrm{ns}}
\newcommand{\h}{\mathrm{h}}
\newcommand{\yr}{\mathrm{yr}}
\newcommand{\N}{\mathrm{N}}
\newcommand{\J}{\mathrm{J}}
\newcommand{\eV}{\mathrm{eV}}
\newcommand{\W}{\mathrm{W}}
\newcommand{\Pa}{\mathrm{Pa}}

% derived units
\newcommand{\mps}{\m\,\s^{-1}}
\newcommand{\mph}{\mi\,\h^{-1}}
\newcommand{\mpss}{\m\,\s^{-2}}

% draft comments
\newcommand{\hoggsay}[1]{{\color{blue}\textbf{#1}}}

\begin{document}

\startps{7}
Due Friday 2026 October 23 at 18:00 on Brightspace.
Be sure to explicitly give credit where it is due, on all problems.
Also credit buddies, web pages, and LLMs (as per course policy).
Failure to give explicit credits will result in points deducted.

\begin{problem}
  For this problem you will write a simple Euler's method code, and write a simple leapfrog code.
  We're going to test these two codes on a very simple problem, the two-body problem in gravity
  (the Kepler problem).
  \subpart Look up the mass of the Sun and Jupiter, and the distance to Jupiter.
  We are going to ignore all the other planets (not a good approximation).
  Compute the circular velocity $v_\mathrm{c}$ at Jupiter's orbit.
  We are going to work in units of years, a.u., and Solar Masses for this problem.
  Put all your answers for this problem---and write your code---in these units.
  \subpart
  Write a function called \code{accelerations(xs)} that takes in six coordinates (the $(x,y,z)$ for the Sun and $(x,y,z)$ for Jupiter and computes six components of the acceleration (the $(a_x, a_y, a_z)$ for the force acting on the Sun and the $(a_x, a_y, a_z)$ acting on Jupiter.
  \subpart
  Start with the Sun at the origin and Jupiter at a reasonable displacement from the Sun (on the $x$ axis), give the Sun an initial velocity of zero and Jupiter an initial velocity (in the $y$ direction) of $0.9\,v_\mathrm{c}$.
  Using both Euler's method and leapfrog, integrate the orbits of the Sun and Jupiter for 100 years.
  Use the same time step for both (choose it wisely; maybe 0.01 years).
  Plot the Euler and leapfrog orbits for Jupiter on a well-labeled plot of $x$ \foreign{vs} $y$.
  Plot the Euler and leapfrog orbits for the Sun on a well-labeled plot, maybe with different axis limits.
  Now zoom in on part of Jupiter's orbit; does the orbit ``close''? At what level does it miss?
  \subpart
  Write a function called \code{angular_momentum(ms, xs, vs)} that takes as input the two masses, the six components of position, and the six components of velocity, and returns the total angular momentum vector about the origin.
  Plot the $z$ component $L_z$ of the angular momentum vs time (in years) for the whole system, for both the Euler and leapfrog orbits.
  Do you have any comments?
\end{problem}

\begin{note}
  To do the last part, you might need to modify your integrators so they save the angular momentum vector as a function of timestep.
\end{note}

\begin{note}
  When you write your leapfrog integrator, you might think a bit about code structure.
  The key things you might want include \code{leapfrog_start()} that moves forward a half-timestep to set up the leapfrog,
  and \code{leapfrog_step()} that moves the system---both positions and velocities---forward by one full step.
  In some applications you might also want \code{leapfrog_end()} that moves things forward by the last half-timestep to get everything back to the same time.
  Finally, if you want good code, make it agnostic to the number of position variables (in this case there are six).
\end{note}

\begin{note}
  If you want a window into Prof Hogg's origins, see \href{https://ui.adsabs.harvard.edu/scan/manifest/1991AJ....101.2274H?art=true}{``Dynamical limits on dark mass in the Solar System'' (1991)}.
  The integrations were for about 100 years, and they were done with a home-built leapfrog integrator.
\end{note}

\begin{problem}
  Newman, \textit{Computational Physics, 2ed}, Exercise~8.3.
\end{problem}

\begin{note}
  You don't have to write your own integrator for this problem.
  You can use any standard Python integrator.
  Or you can write your own, as you wish!
\end{note}

\begin{problem}
  Newman, \textit{Computational Physics, 2ed}, Exercise~8.5.
\end{problem}

\begin{note}
  For this problem you can use your leapfrog integrator from the first \problemname,
  or else you can use any standard Python integrator.
\end{note}

\begin{problem}
  Newman, \textit{Computational Physics, 2ed}, Exercise~8.15.
\end{problem}

\begin{note}
  You don't have to write your own Runge--Kutta integrator for this problem;
  use any standard Python RK4 integrator.
\end{note}

\end{document}
